% Beweis zu FiniteSemigroupTableTupleInjective.
Setze \(S=\NatSeg n\). Die beiden Tupel werden als verschachtelte
endliche Folgen gelesen; die äußere Folge nummeriert die Zeilen.
\begin{tabproofwide}
  \proofstepwidestar[1]{n\in\mathbb N}{\rA}
  \proofstepwidestar[2]{m\colon S\times S\to S}{\rA}
  \proofstepwidestar[3]{d\colon S\times S\to S}{\rA}
  \proofstepwidestar[4]{w_n(m)=w_n(d)}{\rA}
  \proofstepwidestar[5]{i\in S}{\rA}
  \proofstepwidestar[6]{j\in S}{\rA}
  \proofstepwidestarbreak[1,2,3,4,5]{(m(i,j))_{j\in S}=(d(i,j))_{j\in S}}{\FormulaRefAuto{FiniteSemigroupTableTupleDef}[def]{1,2,3},\rIE{4}}
  \proofstepwidestar[]{\bigl((m(i,k))_{k\in S}\bigr)(j)=\bigl((m(i,k))_{k\in S}\bigr)(j)}{\rII}
  \proofstepwidestar[1,2,3,4,5]{\bigl((m(i,k))_{k\in S}\bigr)(j)=\bigl((d(i,k))_{k\in S}\bigr)(j)}{\rIE{7,8}}
  \proofstepwidestar[1,2,3,4,5,6]{m(i,j)=d(i,j)}{\FormulaRefAuto{FiniteSemigroupTableTupleDef}[def]{1,2,3,5,6,9}}
  \proofstepwidestar[1,2,3,4]{\forall i,j\in S\,(m(i,j)=d(i,j))}{\rUI{\rRI{5,\rUI{\rRI{6,10}}}}}
  \proofstepwidestar[1,2,3,4]{m=d}{\FormulaRefAuto{FunctionExtensionality}[thm]{2,3,11}}
\end{tabproofwide}
